\section{拟有限态射与 Zariski 主定理}\label{src:A.7}
\begin{epigraph}
数学是把不同的事物赋予相同名称的艺术.
\epigraphauthor{Henri Poincaré}
\end{epigraph}
概形的局部有限型态射 $f:X\to Y$ 在 $x\in X$ 处\newterm{局部拟有限}{locally quasi-finite}, 如果 $x$ 在纤维 $X_{f(x)}=X\times_Y\Spec\kappa(f(x))$ 中孤立. 当 $f:X\to Y$ 还拟紧时, 这等价于集合 $f^{-1}(f(x))$ 有限, 此时称 $f:X\to Y$ \newterm{拟有限}{quasi-finite}.

\subsection{拟有限态射的平展局部化}\label{src:A.7.1-subsection}
\begin{theorem}[拟有限态射的平展局部化]\label{src:A.7.1}
设 $f:X\to S$ 为分离有限型概形态射. 假设 $f$ 在 $s\in S$ 的每个原像处拟有限. 则存在满足 $\kappa(s')=\kappa(s)$ 的平展邻域 $(S',s')\to(S,s)$, 以及分解 $X\times_S S'=Z\coprod W$, 其中 $Z$ 和 $W$ 为既开又闭的子概形, $Z\to S'$ 有限且纤维 $W_{s'}$ 为空. 此外, 可以使 $Z$ 分解为 $Z_1\coprod\cdots\coprod Z_n$, 其中每个 $Z_i$ 都恰好包含一个位于 $s$ 上方的点 $z_i$, 且 $\kappa(z_i)/\kappa(s)$ 纯不可分.
\end{theorem}
\begin{proof}
见 \cite[IV.$8.12.3$]{EGA} 或 \cite[Tag 04HF]{SP}.
\end{proof}
该结论意味着 Hensel 局部环 $R$ 上的任何拟有限代数 $A$ 都是乘积 $A\simeq B\times C$, 其中 $B$ 在 $R$ 上有限, $C\otimes_R R/\mathfrak m_R=0$; 事实上, 这一性质刻画了 Hensel 局部环 (命题 \bookref{B.5.9}). 它也提供了分解拟有限态射的关键技术.

\subsection{拟有限态射的分解}\label{src:A.7.2-subsection}
\begin{proposition}\label{src:A.7.2}
概形的拟有限分离态射 $f:X\to Y$ 分解为
\[
f:X\longrightarrow\Spec_Y f_*\mathcal O_X\longrightarrow Y,
\]
其中 $X\hookrightarrow\Spec_Y f_*\mathcal O_X$ 为开浸入, $\Spec_Y f_*\mathcal O_X\to Y$ 仿射.
\end{proposition}
\begin{proof}
由于 $f_*\mathcal O_X$ 与 $Y$ 上的平展 (甚至平坦) 基变换交换, 分解也如此. 因而只需证明每个点 $y\in Y$ 都有使命题成立的平展邻域. 由定理 \bookref{A.7.1}, 可假设 $X=X_1\coprod X_2$, 其中 $X_1$ 在 $Y$ 上有限, $(X_2)_y=\varnothing$. 把 $Y$ 换成 $\Spec_Y f_*\mathcal O_X$ 后, 还可以假设 $f_*\mathcal O_X=\mathcal O_Y$. 由于 $\mathcal O_X=\mathcal A_1\times\mathcal A_2$ 是拟凝聚 $\mathcal O_X$-代数的乘积, $\mathcal O_Y=f_*\mathcal O_X=f_*\mathcal A_1\times f_*\mathcal A_2$, 因而 $Y$ 分解为 $Y_1\coprod Y_2$, 满足 $y\in Y_1$ 且对 $i=1,2$, $f(X_i)\subseteq Y_i$. 把 $Y$ 换成 $Y_1$ 后, 可知 $X\to Y$ 有限. 因而 $X$ 仿射, 且 $X=Y=\Spec_Y f_*\mathcal O_X$.
\end{proof}
在上述分解中, $f_*\mathcal O_Y$ 未必为有限型 $\mathcal O_Y$-代数; 即使 $Y$ 为 Noether 仿射概形, $\Gamma(X,\mathcal O_X)$ 也未必为 Noether 环 (见 \cite[例 $7.2.15$]{Ols16}). 然而, 可以修改分解, 使 $X\to Y$ 分解为开浸入后接有限态射.
\begin{theorem}[Zariski 主定理]\label{src:A.7.3}
\end{theorem}
\begin{proof}
若 $\mathcal A\subseteq f_*\mathcal O_X$ 表示 $\mathcal O_Y\to f_*\mathcal O_X$ 的整闭包, 则 $f$ 分解为复合
\[
f:X\xrightarrow{j}\Spec_Y\mathcal A\longrightarrow Y.
\]
断言 $j$ 为开浸入. 只需证明对每个点 $x\in X$, $j(x)$ 都有开邻域 $V\subseteq\Spec_Y\mathcal A$, 使得 $j^{-1}(V)\to V$ 为同构. 由于正规化与平展基变换交换 (命题 \bookref{A.7.4}), 且开浸入是 fpqc 局部性质 (命题 \bookref{3.1.26}), 可以把 $Y$ 换成 $f(x)$ 的平展邻域. 由定理 \bookref{A.7.1}, 可假设 $X=F\coprod W$, 其中 $F$ 在 $Y$ 上有限, $x\in F$. 此时 $Y$ 在 $X$ 中的正规化 $\Spec_Y\mathcal A$ 为 $F\coprod\widetilde W$, 其中 $\widetilde W$ 为 $Y$ 在 $W$ 中的正规化. 由于 $j^{-1}(F)\xrightarrow{\sim}F$, 断言成立. 按构造, $\Spec_Y\mathcal A\to Y$ 为整态射. 可以写 $\mathcal A=\operatorname{colim}\mathcal A_\lambda$, 其中每个 $\mathcal A_\lambda$ 为有限型 $\mathcal O_Y$-代数. 由于开浸入在极限下下降 (命题 \bookref{B.3.7}), 对 $\lambda\gg0$, $X\to\Spec_Y\mathcal A_\lambda$ 为开浸入. 由于 $\Spec_Y\mathcal A_\lambda\to Y$ 为整态射且有限型, 它有限. 另见 \cite[IV.$8.12.6$]{EGA} 或 \cite[Tag 05K0]{SP}.
\end{proof}
以下代数结果已用于上述证明, 也将用于把 Zariski 主定理推广到代数空间 (定理 \bookref{5.5.9}) 和叠 (定理 \bookref{7.1.17}).
\begin{proposition}\label{src:A.7.4}
设 $Y$ 为概形, $\mathcal B$ 为拟凝聚 $\mathcal O_Y$-代数, $\widetilde{\mathcal B}$ 为 $\mathcal O_Y$ 在 $\mathcal B$ 中的整闭包. 若 $f:X\to Y$ 为光滑态射, 则 $f^*\widetilde{\mathcal B}$ 与 $\mathcal O_X$ 在 $f^*\mathcal B$ 中的整闭包等同.
\end{proposition}
\begin{proof}
见 \cite[Tag 03GG]{SP} 或 \cite[命题 $16.2$]{LMB00}.
\end{proof}
Zariski 主定理有一些有用的推论.
\begin{corollary}\label{src:A.7.5}
概形的拟有限固有态射 (相应地, 固有单态射) 有限 (相应地, 为闭浸入).
\end{corollary}
\begin{proof}
若 $f:X\to Y$ 拟有限且固有, 则 Zariski 主定理 (\bookref{A.7.3}) 给出分解 $f:X\hookrightarrow\widetilde X\to Y$, 而稠密开浸入 $X\hookrightarrow\widetilde X$ 也闭, 因而为同构. 另一方面, 若 $f:X\to Y$ 为固有单态射, 则它也拟有限, 因而有限. 结论归结为有限环满态射为满射这一代数事实 (参见 \cite[Tag 04VT]{SP}).
\end{proof}
\begin{remark}\label{src:A.7.6}
由于普遍闭态射必定拟紧 \cite[Tag 04XU]{SP}, 每个普遍闭且局部有限型的单态射都是闭浸入; 另见 \cite[Tag 04XV]{SP}.
\end{remark}
